% Compile with: uplatex file.tex && dvipdfmx file.dvi
%           or: latexmk -uplatex file.tex
\documentclass[dvipdfmx]{article}
\usepackage{amsmath, amssymb}
\usepackage{tikz}

\title{scenario\_coupled\_oscillations}
\date{\today}

\begin{document}
\maketitle

Consider two point masses of mass $m$ connected by two outer springs (constant $k$) and one coupling spring (constant $c$). Let $x,y$ be the displacements of the two masses.

Define x as a time-dependent function

$\checkmark$ Declared x as time-dependent (time variable: t)

Define y as a time-dependent function

$\checkmark$ Declared y as time-dependent (time variable: t)

Define m, k, c as scalars

\[
  m, k, c
\]

Define the Lagrangian $L=$ kinetic energy $-$ potential energy. The potential terms are the outer-spring energies $\frac12kx^2,\ \frac12ky^2$ for each mass, plus the coupling-spring energy $\frac12c(x-y)^2$, proportional to the square of the displacement difference $(x-y)$.

$Lag = (1/2) m \dot{x}^2 + (1/2) m \dot{y}^2 - (1/2) k x^2 - (1/2) k y^2 - (1/2) c (x - y)^2$

\begin{equation}
  Lag = (1/2) m \dot{x}^2 + (1/2) m \dot{y}^2 - (1/2) k x^2 - (1/2) k y^2 - (1/2) c (x - y)^2
  \label{eq:E1}
\end{equation}

Following the Euler-Lagrange equation $\frac{d}{dt}\frac{\partial L}{\partial\dot x}-\frac{\partial L}{\partial x}=0$, derive the equation of motion for $x$: differentiate $L$ with respect to $\dot x$, take the total time derivative of that result, differentiate $L$ with respect to $x$, then take the difference of the two.

Take the partial derivative of E1 with respect to $\dot{x}$

\begin{equation}
  \frac{\partial Lag}{\partial \dot{x}} = m \dot{x}
  \label{eq:E2}
\end{equation}

Take the total derivative of E2 with respect to t

\begin{equation}
  \frac{d}{dt}\left(\frac{\partial Lag}{\partial \dot{x}}\right) = m \ddot{x}
  \label{eq:E3}
\end{equation}

Take the partial derivative of E1 with respect to x

\begin{equation}
  \frac{\partial Lag}{\partial x} = -\left(c \left(x-y\right)\right)-k x
  \label{eq:E4}
\end{equation}

$EOMx = E3 - E4$

\begin{equation}
  EOMx = -\left(c y\right)+m \ddot{x}+\left(k+c\right) x
  \label{eq:E5}
\end{equation}

Derive the equation of motion for $y$ by the same procedure.

Take the partial derivative of E1 with respect to $\dot{y}$

\begin{equation}
  \frac{\partial Lag}{\partial \dot{y}} = m \dot{y}
  \label{eq:E6}
\end{equation}

Take the total derivative of E6 with respect to t

\begin{equation}
  \frac{d}{dt}\left(\frac{\partial Lag}{\partial \dot{y}}\right) = m \ddot{y}
  \label{eq:E7}
\end{equation}

Take the partial derivative of E1 with respect to y

\begin{equation}
  \frac{\partial Lag}{\partial y} = c \left(x-y\right)-k y
  \label{eq:E8}
\end{equation}

$EOMy = E7 - E8$

\begin{equation}
  EOMy = m \ddot{y}+\left(k+c\right) y-c x
  \label{eq:E9}
\end{equation}

Taking the sum and difference of the two equations of motion decouples them into normal modes. The sum cancels the coupling constant $c$, giving the symmetric mode (lower frequency, $\omega_1^2=k/m$); the difference brings in $k+2c$, giving the antisymmetric mode (higher frequency, $\omega_2^2=(k+2c)/m$).

$sym = E5 + E9$

\begin{equation}
  sym = m \ddot{y}+k y+m \ddot{x}+k x
  \label{eq:E10}
\end{equation}

$anti = E5 - E9$

\begin{equation}
  anti = -\left(m \ddot{y}\right)+\left(-k-2 c\right) y+m \ddot{x}+\left(k+2 c\right) x
  \label{eq:E11}
\end{equation}

Define the normal-coordinate transformation $Q_1=x+y$ (the symmetric-mode coordinate) and $Q_2=x-y$ (the antisymmetric-mode coordinate).

$Q_1 = x_1 + x_2$

\begin{equation}
  Q_1 = x_1 + x_2
  \label{eq:E12}
\end{equation}

$Q_2 = x_1 - x_2$

\begin{equation}
  Q_2 = x_1 - x_2
  \label{eq:E13}
\end{equation}

Invert this transformation to express the original displacements $x,y$ (written as $x_1,x_2$ here) in terms of the normal coordinates $Q_1,Q_2$.

Invert E12, E13 for $x_1$, $x_2$

\begin{equation}
  x_{1} = \frac{Q_{2}+Q_{1}}{2}
  \label{eq:E14}
\end{equation}

\begin{equation}
  x_{2} = \frac{Q_{1}-Q_{2}}{2}
  \label{eq:E15}
\end{equation}

\end{document}